\section{External oscillating electric field\label{sec:External-oscillating-electric-field}}

Oscillating orthogonal electric field can bias the electrospinning
process, particularly the radius of the electrospun fibers. Further,
it can be used as a control mechanism to achieve uniform distributions
of polymer filaments at the collector. In JETSPIN it is possible to
simulate external oscillating electric field. Denoted the electric
field

\begin{equation}
\vec{E}=\left(E_{x},\,E_{y},\,E_{z}\right),
\end{equation}
the directive \texttt{external potential type} (see Table \ref{tab:inputlist})
can be used to select the waveform of the oscillation. The user can
select three types of waveform: a rectangular wave oriented along
x-axis; a rectangular wave produced by a RC circuit; a rotating electric
field of a sinusoidal wave. In the first case the electric field is
given by

\begin{equation}
\vec{E}=\left(E_{x}\left\lceil\sin\omega t+\varphi\right\rceil,\,0,\,0\right),
\end{equation}

where the symbol $\left\lceil .\right\rceil$ denotes the ceiling
function. If the waveform is modeled as produced by an RC circuit,
the exponential decay $E_{x}\exp\left(-t/RC\right)$ is added to the
previous equation whenever the field is switched off, and the
exponential increase $E_{x}\left[1-\exp\left(-t/RC\right)\right]$
is added whenever the field is switched on. Here, $RC$ denotes a
decay time (seconds). The last case describes an orthogonal rotating
electric field. Here the electric field takes the form

\begin{equation}
\vec{E}=\left(E_{x},\,E_{y}\cos\omega t+\varphi,\,E_{z}\sin\omega t+\varphi\right).
\end{equation}

Details on the main effects of orthogonal rotating electric fields
are available in Ref. \cite{lauricella2017effects}.
